\section{Native compression of the neuron-gated prior}
\label{sec:ard-native-compression}
This section concerns the native hierarchical prior, without conditioning any weights. At each hidden vertex let $S_j^\ell\sim\operatorname{Gamma}(\kappa/N,1)$, with fixed $0<\kappa\le1$. All edge factors and all biases are independent symmetric Student variables of fixed degrees of freedom $\nu>2$. An outgoing weight from vertex $j$ is $S_j^\ell G_{ij}^{\ell+1}$. The first layer has no gate on its input coordinates. The hidden activation is globally $A$-Lipschitz, and $b=|\phi(0)|$. Biases are present at every hidden layer and at the linear readout. Constants below may depend on this fixed activation and scalar-law choice.

\begin{theorem}[A polynomial-depth circuit dictionary]
\label{thm:ard-native-compression}
For $0<\delta,\tau\le1/2$, put
\[
 L=H+1+\log\frac{d+1}{\delta\tau}.
\]
With prior probability at least $1-\delta$, the raw depth-$H$ network is uniformly $\tau$-approximable on $[-1,1]^d$ by a rational-weight network of the same activation and hidden depth, with
\[
 \text{width}\le CL,\qquad
 \text{hidden vertices}\le CHL,\qquad
 \text{bits per coefficient}\le CHL.
\]
The approximating networks belong to a deterministic finite dictionary, independent of learner width $N$, whose cardinality obeys
\[
 \log|\mathcal D_{H,d,\delta,\tau}|
 \le C\big[dH L^2+H^2L^3\big].
 \label{eq:ard-dictionary-cardinality}
\]
These are polynomial bounds jointly in depth, input dimension, and logarithmic failure probability and precision.
\end{theorem}
\begin{proof}
Write $m_1=\E|G|<\infty$. The gamma density implies, uniformly for $a=\kappa/N\in(0,1]$ and $0<\epsilon\le1$,
\begin{align*}
 N\E[S\mathbf1_{S\le\epsilon}]&\le C\epsilon,\\
 N\Pr(S>\epsilon)&\le C\log(C/\epsilon),\\
 N\Pr(S>B)&\le C e^{-B/2}\qquad(B\ge1).
\end{align*}
Indeed $1/\Gamma(a)\le Ca$, integration below $\epsilon$ gives the first bound, integration over $(\epsilon,1)$ and $(1,\infty)$ gives the second, and the exponential gamma tail gives the third.

Delete all outgoing columns whose gates are at most $\epsilon$, keeping the same factors and biases elsewhere. Let $z^{(\epsilon)}$ denote the resulting raw function. The gate $S_j^\ell$ is independent of its own feature $h_j^\ell$, because that feature depends only on earlier gates. Thus the usual path majorants give
\[
 P_1\le b+A m_1(d+1),\qquad
 P_\ell\le b+A m_1+A\kappa m_1P_{\ell-1},
\]
where $P_\ell$ bounds the expected uniform magnitude at any native or truncated neuron. For a fixed $K>1$, these imply $P_\ell\le C(d+1)K^\ell$. Subtract the coupled networks one layer at a time. A retained column propagates a previous feature error with expected total absolute gain at most $\kappa m_1$. The deleted-column contribution has expected absolute gain at most $C\epsilon$. Independence from the preceding features therefore gives
\[
 \E\|z-z^{(\epsilon)}\|_\infty
 \le C(d+1)\epsilon(H+1)K^{H+1}.
\]
Choose
\[
 \epsilon=\frac{\delta\tau}{C(d+1)(H+1)K^{H+1}}
\]
with a sufficiently large constant. The raw truncation error exceeds $\tau/2$ with probability at most $\delta/4$. Also $\log(C/\epsilon)\le CL$.

In one hidden layer, the number of retained gates is binomial with mean at most $C\log(C/\epsilon)\le CL$. Its exponential moment is bounded by that of a Poisson variable with the same mean. Since $\log(H/\delta)\le CL$, a Chernoff bound and union over layers show that every layer retains at most $U=\lceil CL\rceil$ vertices, except on an event of probability at most $\delta/4$. Similarly, the maximum of all $HN$ gates is at most
\[
 B_S=C\log(CH/\delta)
\]
except on an event of probability at most $\delta/8$.

After deletion, only the retained vertices can influence later layers. Their induced subnetwork uses at most
\[
 M=(d+1)U+(H-1)U(U+1)+U+1
       \le C(dL+HL^2)
\]
edge factors and biases. The retained index sets depend only on gates and are independent of every factor and bias. The Student tail bound therefore permits a conditional union bound over these at most $M$ variables. Except on an event of probability at most $\delta/8$, all their magnitudes are at most
\[
 B_G=1+(CM/\delta)^{1/\nu}.
\]
Every actual coefficient of the induced subnetwork is then bounded by $B_*=\max\{1,B_S\}B_G$, and $\log(B_*+1)\le CL$. This includes its readout bias. Empty retained layers are allowed; subsequent retained neurons may then be constants due to their biases.

Relabel retained vertices consecutively in each layer; their original learner indices no longer occur. Round every coefficient to a common dyadic mesh of step $q\le1$. A deterministic layerwise Lipschitz estimate bounds the resulting uniform output error by
\[
 Cq(H+1)(d+U+1)K_q^{H+1},
 \qquad K_q=1+b+A(d+U+1)(B_*+1).
\]
Choosing a dyadic $q$ small enough to make this at most $\tau/2$ requires at most $CHL$ bits per coefficient, since $\log K_q\le CL$. The total failure probability used above is at most $3\delta/4$, and the truncation and rounding errors sum to at most $\tau$.

Finally enumerate all layer widths between zero and $U$, all allowed dyadic coefficients, and all edges in their consecutive dense layer blocks. There are at most $M$ coefficients, each with $CHL$ bits, plus at most $CH\log(U+1)$ bits for the layer widths. This gives \eqref{eq:ard-dictionary-cardinality}. The enumeration contains no original learner-neuron addresses and is independent of $N$.
\end{proof}

\begin{corollary}[An actual circuit-complexity converse]
\label{cor:ard-prior-complexity}
Let $\mathcal X\subseteq[-1,1]^d$ and suppose
\[
 p=\mu_N\{\|z-f\|_{\infty,\mathcal X}\le\tau\}>0,
 \qquad0<\tau\le1/2.
\]
Then $f$ has a uniform $2\tau$ approximation on $\mathcal X$ by a rational-weight network of hidden depth $H$, width at most $CL_p$ and coefficient bit length at most $CHL_p$, where
\[
 L_p=H+1+\log\frac{d+1}{p\tau}.
\]
Its description has at most $C[dH L_p^2+H^2L_p^3]$ bits.
\end{corollary}
\begin{proof}
Apply Theorem~\ref{thm:ard-native-compression} with failure probability $p/2$. Its compression event intersects the specified function neighborhood. The triangle inequality proves the approximation claim, and the remaining conclusions are the dictionary bounds.
\end{proof}
